\chapter{Chapter06 — Restriction Mean Variance}
%%% generated by the blueprint pipeline from FABL.Chapter06.Pseudorandomness.RestrictionMeanVariance
\section{Overview}

Declarations of the Lean module \texttt{FABL.Chapter06.Pseudorandomness.RestrictionMeanVariance}.

\section{Declarations}

\begin{theorem}[Variance and nonconstant Fourier coefficients]
\lean{FABL.variance_eq_sum_sq_indexedFourierCoeff_ne_empty}
\proofsource{Analysis-of-Boolean-Functions-by-Ryan-ODonnell}{
  pdf-ffbe9c7b3306-p027-b001,
  pdf-ffbe9c7b3306-p027-b002,
  pdf-ffbe9c7b3306-p211-b005
}
\label{FABL.variance_eq_sum_sq_indexedFourierCoeff_ne_empty}
\leanok
\uses{FABL.IndexedSignCube, FABL.expect_eq_indexedFourierCoeff_empty, FABL.indexedFourierCoeff, FABL.indexed_parseval, FABL.mean, FABL.variance}
\difficulty{4}
Let $\iota$ be a finite index set and let $g$ be a real-valued function on the sign cube
$\{\pm1\}^{\iota}$. Its variance under the uniform distribution is the sum, over all
nonempty subsets $S\subseteq\iota$, of the squares of its Fourier coefficients indexed
by $S$.
\end{theorem}

\begin{theorem}[Mean of a restriction and its constant coefficient]
\lean{FABL.mean_signRestriction_eq_restrictionFourierCoeff_empty}
\label{FABL.mean_signRestriction_eq_restrictionFourierCoeff_empty}
\leanok
\uses{FABL.FixedSignCube, FABL.expect_eq_indexedFourierCoeff_empty, FABL.indexedFourierCoeff, FABL.mean, FABL.restrictionFourierCoeff, FABL.signRestriction}
\difficulty{1}
Let $J$ be a set of coordinates in the $n$-dimensional sign cube, and fix a sign
assignment $z$ on the complementary coordinates. The mean of the restriction of a
real-valued function $f$ to the free coordinates equals the Fourier coefficient indexed
by the empty set of that restricted function; equivalently, it equals the mean of $f$
over the free coordinates with $z$ held fixed.
\end{theorem}

\begin{theorem}[Variance of restricted means]
\lean{FABL.variance_mean_signRestriction_eq_sum_sq_fixed}
\proofsource{Analysis-of-Boolean-Functions-by-Ryan-ODonnell}{
  pdf-ffbe9c7b3306-p077-b001,
  pdf-ffbe9c7b3306-p077-b002
}
\label{FABL.variance_mean_signRestriction_eq_sum_sq_fixed}
\leanok
\uses{FABL.FixedIndex, FABL.FixedSignCube, FABL.fourierCoeff, FABL.indexedFourierCoeff_restrictionFourierCoeff, FABL.liftFixedFrequency, FABL.liftFreeFrequency, FABL.mean, FABL.mean_signRestriction_eq_restrictionFourierCoeff_empty, FABL.restrictionFourierCoeff, FABL.signRestriction, FABL.variance, FABL.variance_eq_sum_sq_indexedFourierCoeff_ne_empty}
\difficulty{3}
Let $n$ be a natural number, let $J$ be a set of coordinates in $\{0,\ldots,n-1\}$, and
let $f:\{−1,1\}^n\to\mathbb R$. For each assignment of signs to the coordinates outside
$J$, average $f$ over all sign assignments on $J$. The variance of these averages, over
uniformly chosen assignments outside $J$, equals the sum of the squares of the Fourier
coefficients of $f$ indexed by nonempty subsets of coordinates outside $J$.
\end{theorem}

\begin{theorem}[Mean change under a sign restriction]
\lean{FABL.exists_signRestriction_mean_change_gt_via_variance}
\proofsource{Analysis-of-Boolean-Functions-by-Ryan-ODonnell}{
  pdf-ffbe9c7b3306-p146-b006,
  pdf-ffbe9c7b3306-p147-b003
}
\label{FABL.exists_signRestriction_mean_change_gt_via_variance}
\leanok
\uses{FABL.FixedIndex, FABL.FixedSignCube, FABL.IsLowDegreeFourierRegular, FABL.expect_restrictionFourierCoeff, FABL.fixedFrequencyPart, FABL.fourierCoeff, FABL.freeFrequencyPart, FABL.liftFixedFrequency, FABL.liftFreeFrequency, FABL.liftFreeFrequencyPart_union_liftFixedFrequencyPart, FABL.mean, FABL.mean_eq_fourierCoeff_empty, FABL.mean_signRestriction_eq_restrictionFourierCoeff_empty, FABL.mem_freeFrequencyPart, FABL.restrictionFourierCoeff, FABL.signRestriction, FABL.variance, FABL.variance_mean_signRestriction_eq_sum_sq_fixed}
\difficulty{5}
Let $f$ be a real-valued function on the $n$-dimensional sign cube, and suppose it is
not Fourier-regular at degree at most $k$ with tolerance $\varepsilon$; that is, some
nonempty set of coordinates of size at most $k$ has a Fourier coefficient of absolute
value greater than $\varepsilon$. Then there is a set $J$ of coordinates whose
complement has size at most $k$, and an assignment of signs to the coordinates outside
$J$, such that restricting $f$ to that assignment changes its mean by more than
$\varepsilon$ in absolute value.
\end{theorem}
